% interactapasample.tex
% v1.05 - August 2017

\documentclass[]{interact}

\usepackage{epstopdf}
\usepackage[caption=false]{subfig}

\usepackage[longnamesfirst,sort]{natbib}
\bibpunct[, ]{(}{)}{;}{a}{,}{,}
\renewcommand\bibfont{\fontsize{10}{12}\selectfont}

\theoremstyle{plain}
\newtheorem{theorem}{Theorem}[section]
\newtheorem{lemma}[theorem]{Lemma}
\newtheorem{corollary}[theorem]{Corollary}
\newtheorem{proposition}[theorem]{Proposition}

\theoremstyle{definition}
\newtheorem{definition}[theorem]{Definition}
\newtheorem{example}[theorem]{Example}

\theoremstyle{remark}
\newtheorem{remark}{Remark}
\newtheorem{notation}{Notation}

\begin{document}

\articletype{ARTICLE TEMPLATE}

\title{Taylor \& Francis \LaTeX\ template for authors (\textsf{Interact} layout + American Psychological Association reference style)}

\author{
\name{A.~N. Author\textsuperscript{a}\thanks{CONTACT A.~N. Author. Email: latex.helpdesk@tandf.co.uk} and John Smith\textsuperscript{b}}
\affil{\textsuperscript{a}Taylor \& Francis, 4 Park Square, Milton Park, Abingdon, UK; \textsuperscript{b}Institut f\"{u}r Informatik, Albert-Ludwigs-Universit\"{a}t, Freiburg, Germany}
}

\maketitle

\begin{abstract}
This template is for authors who are preparing a manuscript for a Taylor \& Francis journal using the \LaTeX\ document preparation system and the \textsf{interact} class file, which is available via selected journals' home pages on the Taylor \& Francis website.
\end{abstract}

\begin{keywords}
Sections; lists; figures; tables; mathematics; fonts; references; appendices
\end{keywords}

\section{Introduction}

In order to assist authors in the process of preparing a manuscript for a journal, the Taylor \& Francis `Interact' layout style has been implemented as a \LaTeXe\ class file based on the \textsf{article} document class. A sample bibliography is also provided in order to assist with the formatting of your references.

Commands that differ from or are provided in addition to standard \LaTeXe\ are described in this document, which is \textit{not} a substitute for a \LaTeXe\ tutorial.

The \texttt{interactapasample.tex} file can be used as a template for a manuscript by cutting, pasting, inserting and deleting text as appropriate, using the preamble and the \LaTeX\ environments provided (e.g.\ \texttt{\textbackslash begin\{abstract\}}, \texttt{\textbackslash begin\{keywords\}}).

\subsection{The \textsf{Interact} class file}

The \textsf{interact} class file preserves the standard \LaTeXe\ interface such that any document that can be produced using \textsf{article.cls} can also be produced with minimal alteration using the \textsf{interact} class file as described in this document.

\section{Some {\itshape Lancet} style headings}

\subsection{This is a level-2 heading}

This is normal paragraph text.  This is normal paragraph text.  This is normal paragraph text.  This is normal paragraph text.

\subsubsection{This is a level-3 heading}

This is normal paragraph text.  This is normal paragraph text.  This is normal paragraph text.

\paragraph{This is a level-4 heading}

This is normal paragraph text.  This is normal paragraph text.

\subparagraph{This is a level-5 heading}

This is normal paragraph text.

\section{Notes on the sample files and production of your manuscript}

Please supply any additional macro definitions you use with your article in the preamble of your \texttt{.tex} file.

\section{The references}

The list of references is headed `References' and to assist in the formatting of the references a sample bibliography has been provided with this template.

\bibliographystyle{apalike}
\bibliography{interactapasample}

\end{document}
